\begin{enumerate}[itemsep=5mm]
	\item On a
	\[
	f(x)=4\e^{-x}+\cos x+\sin x.
	\]
	Donc :
	\[
	f'(x)=-4\e^{-x}-\sin x+\cos x.
	\]
	Ainsi :
	\[
	f'(x)+f(x)=\left(-4\e^{-x}-\sin x+\cos x\right)+\left(4\e^{-x}+\cos x+\sin x\right)=2\cos x.
	\]
	La fonction $f$ est bien solution de l'équation différentielle $(E)$.

	\medskip
	\textbf{Affirmation 1 : vraie.}

	\item Les points d'intersection des courbes vérifient :
	\[
	2x=\sin x.
	\]
	Posons $h(x)=2x-\sin x$.
	Alors :
	\[
	h'(x)=2-\cos x.
	\]
	Or, pour tout réel $x$, $\cos x\leqslant 1$, donc :
	\[
	h'(x)=2-\cos x\geqslant 1>0.
	\]
	La fonction $h$ est donc strictement croissante sur $\R$.
	De plus :
	\[
	h(0)=0.
	\]
	L'équation $h(x)=0$ admet donc une unique solution, à savoir $x=0$.
	Les deux courbes ont donc un seul point d'intersection, le point de coordonnées $(0~;~0)$.

	\medskip
	\textbf{Affirmation 2 : vraie.}

	\item Pour tout entier naturel $n$ :
	\[
	v_n=\dfrac{2n+\sin n}{n+1}.
	\]
	On écrit :
	\[
	v_n=\dfrac{n\left(2+\frac{\sin n}{n}\right)}{n\left(1+\frac{1}{n}\right)}
	=\dfrac{2+\frac{\sin n}{n}}{1+\frac{1}{n}}.
	\]
	Comme $-1\leqslant \sin n\leqslant 1$, on a
	\[
	\lim\limits_{n\to +\infty}\dfrac{\sin n}{n}=0.
	\]
	Donc :
	\[
	\lim\limits_{n\to +\infty} v_n=\dfrac{2+0}{1+0}=2.
	\]
	La suite $(v_n)$ converge vers 2.

	\medskip
	\textbf{Affirmation 3 : fausse.}

	\item Montrons par récurrence que, pour tout entier naturel $n\geqslant 1$, $u_n=n^2$.

	\emph{Initialisation :} pour $n=1$, on a $u_1=1$ et $1^2=1$. La propriété est vraie au rang 1.

	\emph{Hérédité :} supposons que, pour un entier $n\geqslant 1$, on ait $u_n=n^2$.
	Alors :
	\[
	u_{n+1}=u_n+2n+1=n^2+2n+1=(n+1)^2.
	\]
	La propriété est donc héréditaire.

	Par récurrence, pour tout entier naturel $n\geqslant 1$, on a $u_n=n^2$.

	\medskip
	\textbf{Affirmation 4 : vraie.}

	\item La suite $(u_n)$ est géométrique de premier terme $u_0=1$ et de raison $\e^{-1}$.
	Comme $0<\e^{-1}<1$, on a :
	\[
	\lim\limits_{n\to +\infty} S_n=\dfrac{u_0}{1-\e^{-1}}
	=\dfrac{1}{1-\frac{1}{\e}}
	=\dfrac{\e}{\e-1}.
	\]

	\medskip
	\textbf{Affirmation 5 : vraie.}
\end{enumerate}
